% from ega1-4-fr.tex
\oldpage[I]{123}
  directe $(\psi_V)_*(\mathscr{O}_Y|V)$, où $\psi_V$ est la restriction de
  $\psi$ à $V$. Or, la restriction à $(V,\mathscr{O}_Y|V)$ du morphisme
  $(\psi,\theta)$ est un morphisme de ce préschéma dans
  $(X,\mathscr{O}_X)$, et par suite est de la forme
  $({}^a\varphi,\widetilde{\varphi})$, où $\varphi$ est un homomorphisme de
  l'anneau $A=\Gamma(X,\mathscr{O}_X)$ dans l'anneau
  $\Gamma(V,\mathscr{O}_Y)$ (\hyperref[I.1.7.3-fr]{1.7.3})~; on en conclut
  que $(\psi_V)_*(\mathscr{O}_Y|V)$ est un $\mathscr{O}_X$-Module
  quasi-cohérent (\hyperref[I.1.6.3-fr]{1.6.3}), ce qui prouve notre
  assertion, en raison du caractère local des faisceaux quasi-cohérents. En
  outre, l'hypothèse que $\psi$ est un homéomorphisme entraîne
  (\hyperref[0.3.4.5-fr]{0, 3.4.5}) que pour tout $y\in Y$, $\psi_y$ est un
  isomorphisme
  $(\psi_*(\mathscr{O}_Y))_{\psi(y)}\to\mathscr{O}_y$~; comme le diagramme
  \phantomsection\label{I.4.2.2.theta-diagram-fr}
  \[
    \xymatrix{
      \mathscr{O}_{\psi(y)}\ar[r]^{\theta_{\psi(y)}}
        \ar[d]_{\psi_y\circ\rho_{\psi(y)}} &
      (\psi_*(\mathscr{O}_Y))_{\psi(y)}\ar[d]^{\psi_y}\\
      (\psi^*(\mathscr{O}_X))_y\ar[r]^{\theta_y^\sharp} &
      \mathscr{O}_y
    }
  \]
  est commutatif et que les deux flèches verticales sont des isomorphismes
  (\hyperref[0.3.7.2-fr]{0, 3.7.2}), l'hypothèse que $\theta_y^\sharp$ est
  surjectif entraîne qu'il en est de même de $\theta_{\psi(y)}$. Comme le
  support de $\psi_*(\mathscr{O}_Y)$ est $Z=\psi(Y)$, $\theta$ est un
  homomorphisme \emph{surjectif} de $\mathscr{O}_X=\widetilde{A}$ dans le
  $\mathscr{O}_X$-Module quasi cohérent $f_*(\mathscr{O}_Y)$. Par suite, il
  existe un isomorphisme unique $\omega$ d'un faisceau quotient
  $\widetilde{A}/\widetilde{\mathfrak{J}}$ ($\mathfrak{J}$ idéal de $A$) sur
  $f_*(\mathscr{O}_Y)$ qui, composé avec l'homomorphisme canonique
  $\widetilde{A}\to\widetilde{A}/\widetilde{\mathfrak{J}}$, donne $\theta$
  (\hyperref[I.1.3.8-fr]{1.3.8})~; si $\mathscr{O}_Z$ désigne la restriction
  de $\widetilde{A}/\widetilde{\mathfrak{J}}$ à $Z$,
  $(Z,\mathscr{O}_Z)$ est un sous-préschéma de $(X,\mathscr{O}_X)$, et $f$ se
  factorise en l'injection canonique de ce sous-préschéma dans $X$, et
  l'isomorphisme $(\psi_0,\omega_0)$, où $\psi_0$ est $\psi$ considéré comme
  application de $Y$ sur $Z$, et $\omega_0$ la restriction de $\omega$ à
  $\mathscr{O}_Z$.

  Passons au cas général. Soit $U$ un ensemble ouvert affine dans $X$ tel que
  $U\cap\psi(Y)$ soit fermé dans $U$ et non vide. En restreignant $f$ au
  préschéma induit par $Y$ sur l'ouvert $\psi^{-1}(U)$, et en le considérant
  comme un morphisme de ce préschéma dans le préschéma induit par $X$ sur
  $U$, on est ramené au premier cas~; la restriction de $f$ à
  $\psi^{-1}(U)$ est donc une immersion fermée $\psi^{-1}(U)\to U$, se
  factorisant canoniquement en $j_U\circ g_U$, où $g_U$ est un isomorphisme
  du préschéma $\psi^{-1}(U)$ sur un sous-préschéma $Z_U$ de $U$, et $j_U$
  l'injection canonique $Z_U\to U$. Soit $V$ un second ouvert affine de $X$
  tel que $V\subset U$~; comme la restriction $Z'_V$ de $Z_U$ à $V$ est un
  sous-préschéma du préschéma $V$, la restriction de $f$ à $\psi^{-1}(V)$ se
  factorise en $j'_V\circ g'_V$, où $j'_V$ est l'injection canonique
  $Z'_V\to V$ et $g'_V$ un isomorphisme de $\psi^{-1}(V)$ sur $Z'_V$. Par
  l'unicité de la factorisation canonique d'une immersion
  (\hyperref[I.4.2.1-fr]{4.2.1}), on a nécessairement $Z'_V=Z_V$ et
  $g'_V=g_V$. On en conclut (\hyperref[I.4.1.5-fr]{4.1.5}) qu'il y a un
  sous-préschéma $Z$ de $X$ dont l'espace sous-jacent est $\psi(Y)$ et dont
  la restriction à chaque $U\cap\psi(Y)$ est $Z_U$~; les $g_U$ sont alors les
  restrictions aux $\psi^{-1}(U)$ d'un isomorphisme $g:Y\to Z$ tel que
  $f=j\circ g$, où $j$ est l'injection canonique $Z\to X$.
\end{enumerate}

\begin{corollary}[4.2.3]
\label{I.4.2.3-fr}
Soit $X$ un schéma affine. Pour qu'un morphisme
$f=(\psi,\theta):Y\to X$ soit une immersion fermée, il faut et il suffit que
$Y$ soit un schéma affine et que l'homomorphisme
$\Gamma(\psi):\Gamma(\mathscr{O}_X)\to\Gamma(\mathscr{O}_Y)$ soit surjectif.
\end{corollary}

\begin{corollary}[4.2.4]
\label{I.4.2.4-fr}
\begin{enumerate}
  \item[a)] Soient $f$ un morphisme $Y\to X$, $(V_\lambda)$ un recouvrement
  de $f(Y)$ par des ouverts de $X$. Pour que $f$ soit une immersion (resp. une
  immersion ouverte), il faut et il suffit
